\documentclass[11pt]{article}
\usepackage[margin=1in]{geometry}
\usepackage{amsmath,amssymb,amsthm,booktabs,graphicx,xcolor,hyperref,enumitem}
\setlist{itemsep=2pt,topsep=3pt}
\newcommand{\passk}{\text{pass@}k}
\newtheorem{proposition}{Proposition}
\newtheorem{assumption}{Assumption}
\theoremstyle{definition}
\newtheorem{definition}{Definition}

\title{\textbf{Coverage Preservation in RLVR as an Off-Policy Constraint}\\
\large A base-anchored support floor for the reasoning coverage GRPO cannot protect}
\author{Methodology and initial results \\ \small Qwen2.5-Math (1.5B / 7B) $\cdot$ MATH-500, OlympiadBench, AIME}
\date{}

\begin{document}
\maketitle

\begin{abstract}
RLVR (e.g.\ GRPO) improves an LLM's single-sample accuracy (pass@1) but is observed to \emph{narrow}
its reasoning \emph{coverage}: at large sampling budget $k$, the base model solves problems the
RLVR model no longer can \citep{yue2025}. We frame the relevant resource not as entropy but as
\emph{finite-budget recoverability}, and make three contributions. \textbf{(1) A mechanism result:}
the expected on-policy gradient routed to a correct reasoning mode $M$ scales as $O(p_\theta(M))$, so
any estimator whose gradient support is a finite set of $K$ on-policy samples becomes
\emph{support-blind} to modes of mass $p_\theta(M)\!\ll\!1/K$ --- no reward/advantage reshaping
(UCPO, PKPO, RiskPO) escapes the factor. \textbf{(2) A new object and constraint:} we preserve the
\emph{recoverability} $R_K(M)=1-(1-p_\theta(M))^K$ of verified reasoning modes with an
\emph{off-policy, one-sided, base-anchored} support floor $p_\theta(M)\ge\alpha\,p_0(M)$, enforced by
teacher forcing on a bank of base-correct traces; its gradient does not vanish as $p_\theta\to0$, and
a Jensen bound turns the banked traces into a \emph{finite-compute recoverability certificate}. Two
implementations: soft Lagrangian (\texttt{expSR}) and hard projected-gradient (\texttt{expPROJ}).
\textbf{(3) Evidence.} The crossover is a difficulty \emph{resonance} (Oat-Zero-7B: present on
intermediate Olympiad, absent when saturated/too-hard). At the reasoning-mode level (held-out
11{,}664-witness bank), identical continued RL from a plain fork \emph{collapses} $39\%$ of base
modes ($\Delta<\!-10$ nats) while a floor fork collapses $2.3\%$; an \emph{interventional} run
(identical GRPO from base, floor the sole difference) is causal --- the floor raises mean mode mass
($+1.38$ nats, $0\%$ collapse) where plain RL erodes it ($-1.96$, $4.1\%$), and it beats global-KL
and PBA/base-anchoring baselines under matched training (it is the only arm with positive mean $\Delta$
and zero collapse). Finally, coverage is the \emph{substrate for continued RL}: a floor-constrained
round-1 fork, after \emph{identical} round-2 RL, \emph{overtakes} the plain fork at large $k$
($k\!\gtrsim\!32$; $0.990$ pass@256) while the plain fork wins small $k$ --- the round-1 pass@1 tax is
recovered as a higher round-2 coverage ceiling. We report honest limits: gains on the saturated
fragile band are small ($\sim\!0.7$ pt); the multi-point, multi-model, multi-seed campaign is future work.
\end{abstract}

%====================================================================
\section{Problem setup}
Let $\pi_\theta$ be the policy LLM and $\pi_0$ the frozen base model. For a prompt $q$, RLVR maximizes
the expected verifiable reward $J(\theta)=\mathbb{E}_{q}\mathbb{E}_{y\sim\pi_\theta(\cdot\mid q)}[r(q,y)]$,
$r\in\{0,1\}$. We measure the \emph{reasoning coverage} of a model by
$\passk(q)=\Pr[\exists\text{ correct among }k\text{ i.i.d.\ samples}]$, averaged over $q$. The
empirical phenomenon \citep{yue2025} is a \emph{crossover}: RLVR raises $\passk$ at small $k$ but base
$\ge$ RLVR at large $k$, i.e.\ RLVR trades coverage for sampling efficiency.

%====================================================================
\section{Mechanism: why reward-shaping cannot fix coverage}
\label{sec:mech}
Consider GRPO with $K$ on-policy rollouts per prompt and positive normalized advantage $A^{+}$ on
correct rollouts. For a specific correct reasoning mode $y$ (a set of trajectories reaching a correct
answer along one solution strategy), the expected policy-gradient contribution routed to $y$ is
\begin{equation}
\mathbb{E}\big[g_y\big]
= A^{+}\cdot K\cdot \pi_\theta(y\mid q)\,\nabla_\theta\log\pi_\theta(y\mid q),
\qquad\text{so}\qquad
\big\lVert \mathbb{E}[g_y]\big\rVert = O\!\big(\pi_\theta(y\mid q)\big).
\label{eq:vanish}
\end{equation}
The factor $\pi_\theta(y\mid q)$ appears because $y$ contributes to the objective only when it is
\emph{sampled}, and it is sampled with probability $\pi_\theta(y\mid q)$.

\begin{assumption}[Sharpening]\label{ass:sharp}
During RLVR the policy concentrates on its dominant correct mode(s); a rare-correct mode $y$ with
$\pi_\theta(y\mid q)\ll 1/K$ is sampled with probability $\to 0$ over training.
\end{assumption}

\begin{proposition}[Irreversible pruning under on-policy shaping]\label{prop:prune}
Under Assumption~\ref{ass:sharp}, for any method whose update modifies only the advantage/reward
weighting $A^{+}\!\mapsto\! w\,A^{+}$ with $w$ bounded, the expected gradient on a mode with
$\pi_\theta(y\mid q)\to 0$ satisfies $\lVert\mathbb{E}[g_y]\rVert\to 0$. Hence once a rare-correct mode
drops below the sampling threshold it receives asymptotically no learning signal and cannot be
restored by any such method.
\end{proposition}

\noindent\emph{Proof sketch.} Reweighting $A^{+}$ by a bounded $w$ leaves the $O(\pi_\theta)$ factor
in \eqref{eq:vanish} intact; as $\pi_\theta(y\mid q)\to0$ the product $\to0$. UCPO's uniformity
penalty, PKPO's pass@$k$ reward transform, RiskPO's risk weighting, and Polychromic's diversity
advantage all instantiate a bounded $w$. $\qed$

\medskip
\noindent This is the gap we target: the modes RLVR loses are exactly those it has stopped sampling,
and no on-policy reward can push on a mode it does not sample.

%====================================================================
\section{Method: an off-policy base-anchored support floor}
\label{sec:method}
\paragraph{Coverage bank.} Offline, we run the base model on the difficulty-labeled \emph{fragile
band} (problems with low but nonzero base pass@1) and collect a bank
$\mathcal{B}=\{(q,y_q)\}$ of base-\emph{correct} traces, storing each trace's reference log-probability
$\ell^{\mathrm{ref}}_q=\log\pi_0(y_q\mid q)$.

\begin{definition}[Support floor]
For $\alpha\in(0,1]$, the policy \emph{satisfies the floor} on $(q,y_q)$ iff
\begin{equation}
\log\pi_\theta(y_q\mid q)\;\ge\;\ell^{\mathrm{ref}}_q+\log\alpha
\qquad(\text{$\alpha{=}0.5$: policy mass may not fall below half the base's on that trace}).
\label{eq:floor}
\end{equation}
\end{definition}

Three properties distinguish this from all reward-shaping. \textbf{(i) Off-policy / teacher-forced:}
$\log\pi_\theta(y_q\mid q)$ is evaluated on the \emph{banked} trace, so the correction gradient's
magnitude is independent of the current $\pi_\theta(y_q\mid q)$. \textbf{(ii) A constraint, not a
reward:} it acts on $\log\pi_\theta$ directly, immune to the ``constant-across-group $\Rightarrow$
no-op'' cancellation in GRPO's group-relative advantage. \textbf{(iii) One-sided (a ratchet):} only
dropping \emph{below} the floor is penalized, so the constraint never pulls the policy back toward
base and imposes no intrinsic pass@1 tax.

\begin{proposition}[Non-vanishing rescue]\label{prop:rescue}
Let $g^{\mathrm{proj}}_y=\nabla_\theta\log\pi_\theta(y_q\mid q)$ be the (teacher-forced) constraint
gradient and $g^{\mathrm{RL}}_y$ the on-policy reward gradient of \eqref{eq:vanish}. In the pruning
regime $\pi_\theta(y_q\mid q)\to0$ (Assumption~\ref{ass:sharp}),
$\lVert\mathbb{E}[g^{\mathrm{RL}}_y]\rVert / \lVert g^{\mathrm{proj}}_y\rVert \to 0$: the constraint
delivers $\Omega(1)$ signal exactly where the reward delivers $o(1)$.
\end{proposition}

\paragraph{Two implementations.}
\emph{Soft (\texttt{expSR}).} Add a one-sided Lagrangian penalty to the GRPO loss:
$\mathcal{L}=\mathcal{L}_{\mathrm{GRPO}}+\mu\,\mathbb{E}_{(q,y_q)\sim\mathcal{B}}
[\mathrm{relu}(\ell^{\mathrm{ref}}_q+\log\alpha-\log\pi_\theta(y_q\mid q))]$, $\mu{=}0.5$; the penalty
is $0$ when feasible.
\emph{Hard (\texttt{expPROJ}).} After each GRPO step, run up to $m_{\max}$ projected-gradient
correction sub-steps that descend the violation on any banked traces below their floor, restoring
feasibility. Cost is bounded by a rotating-window scan (check $32$ traces every $4$ steps,
$m_{\max}{=}2$, correction lr $10^{-5}$; $\sim$60s/step vs.\ $\sim$600s for a full-bank scan).
A \emph{mode-survival} diagnostic logs the fraction of banked base-correct modes still above their
floor during training.

%====================================================================
\section{Results}
\label{sec:results}

\subsection{The phenomenon on a released full-RL model (no self-training confound)}
Our own LoRA/light-trained models compress the crossover (\S\ref{sec:recipe}); the phenomenon
requires \emph{full-parameter, heavily-trained} RL, exactly the released checkpoints
\citet{yue2025} use. We therefore evaluate \textbf{Oat-Zero-7B} (Qwen2.5-Math-7B after full-parameter
zero-RL) against its base, so no confound from our training. This is where the real evidence lives.

\paragraph{(R1) The crossover is a resonance, confirmed across difficulty.}
Base vs Oat-Zero pass@256 (Table~\ref{tab:oatcross}): the crossover appears \emph{only} at
intermediate difficulty. On AMC (easy) both saturate at $1.0$; on \textbf{OlympiadBench
(intermediate)} base overtakes Oat-Zero at $k{\approx}128$ (a \textbf{22-problem lost set}); on AIME (too
hard) Oat-Zero leads everywhere (no samplable base tail). The effect is real but localized to a
\emph{fragile band}. Crucially, the AIME gap is not static: extending to $k{=}1024$
(base vs Oat $0.057/0.165$ at $k{=}1$ $\to$ $0.722/0.733$ at $k{=}1024$) the deficit closes
\emph{monotonically} from $-0.109$ to $-0.011$---the curves are converging toward a crossover just
beyond the measured range, exactly the ``too-hard boundary'' the resonance picture predicts.

\begin{table}[h]\centering\small
\begin{tabular}{llrrr}
\toprule
Benchmark (base vs Oat-Zero-7B) & difficulty & base @256 & Oat @256 & crossover \\
\midrule
AMC ($n{=}40$)          & easy         & 1.000 & 1.000 & none (saturated)\\
\textbf{OlympiadBench ($n{=}572$)} & \textbf{intermediate} & \textbf{0.853} & 0.837 & \textbf{yes: 22-problem lost set}\\
AIME ($n{=}90$)         & hard         & 0.560 & \textbf{0.604} & none (gap $\to0$ by $k{=}1024$)\\
\bottomrule
\end{tabular}
\caption{Crossover is a difficulty resonance, on a released full-RL model. Base overtakes RL only at
intermediate difficulty (Olympiad). See Table~\ref{tab:oly1024} for the $k{=}1024$ curve (widens).}
\label{tab:oatcross}
\end{table}

\begin{table}[h]\centering\small
\begin{tabular}{lrrrrrrr}
\toprule
OlympiadBench ($n{=}572$, full) & pass@1 & pass@64 & pass@128 & pass@256 & pass@512 & pass@1024\\
\midrule
base      & 0.217 & 0.772 & \textbf{0.815} & \textbf{0.853} & \textbf{0.882} & \textbf{0.904}\\
Oat-Zero  & \textbf{0.447} & \textbf{0.775} & 0.808 & 0.837 & 0.864 & 0.888\\
\bottomrule
\end{tabular}
\caption{\textbf{The crossover widens monotonically with $k$} (full $n{=}572$, $1024$ samples/problem).
Base overtakes Oat-Zero at $k{\approx}128$ and the gap grows to $+0.018$ at $k{=}512$, holding
$+0.016$ at $k{=}1024$. Lost-set: $22$ problems at $k{=}1024$. This is the regime where the support
floor has maximal headroom to recover coverage.}
\label{tab:oly1024}
\end{table}

\paragraph{(R2) Mechanism, measured directly --- RL suppresses the exact modes it loses.}
On the $49$ base-correct traces for problems Oat-Zero \emph{lost} on OlympiadBench (base solves @256,
Oat does not): median base $\log\pi = -254 \to$ median Oat-Zero $\log\pi = -320$, and \textbf{$49/49$
traces ($100\%$) are suppressed by Oat-Zero}, mean $-84$ nats.
\textbf{Replication on Omni-MATH ($n{=}150$ base-correct traces, hard band):}
median base $\log\pi=-170\to$ Oat-Zero $-249$, \textbf{$150/150$ ($100\%$) suppressed}, mean $-75$
nats. The mechanism replicates across benchmarks: RLVR does not merely fail on these problems; it
drives \emph{every} recoverable solution $75$--$84$ nats below the base's probability---far under
the sampling floor, which is precisely why they cannot be recovered at large $k$. This is the
empirical realization of Prop.~\ref{prop:prune} and the paper's central mechanism figure.

\paragraph{(R3) Exploitability gate --- the preserved mass is usable by continued RL.}
Pre-flight test (\S\ref{sec:round2}) on Oat-Zero over the $329$-problem Olympiad protect-set at
round-2 sampling ($G{=}16$, $T{=}1.0$): \textbf{$66.3\%$ of protected problems yield a nonzero-advantage
GRPO group} (both a correct and an incorrect rollout), $78.4\%$ produce $\ge1$ correct. The failure
mode ``preserved probability too sparse to exploit'' does \emph{not} occur; the round-2 headline
(\S\ref{sec:round2}) is viable without resorting to the set-form floor or $\alpha>1$.

\subsection{Method ablation on our own models (directional, within-noise at this scale)}
\textbf{Setup.} Qwen2.5-Math, GRPO+LoRA($r{=}32$)+ZeRO-2, $500$ steps; bank on the difficulty-labeled
fragile band, $\alpha{=}0.5$. Evaluation: $\passk$ with $256$--$1024$ samples/problem, $8$-way
data-parallel sharding; paired-bootstrap $95\%$ CIs over problems ($B{=}5000$) vs.\ GRPO.

\begin{table}[h]\centering\small
\begin{tabular}{lrrrrrr}
\toprule
Arm (7B, OlympiadBench, $n{=}572$) & pass@1 & pass@16 & pass@32 & pass@64 & pass@128 & pass@256\\
\midrule
base                & 0.2163 & 0.6449 & 0.7115 & 0.7621 & 0.8017 & 0.8339\\
GRPO (control)      & \textbf{0.2354} & \textbf{0.6565} & 0.7224 & 0.7739 & 0.8154 & 0.8497\\
\texttt{expSR}      & 0.2214 & 0.6482 & 0.7170 & 0.7722 & 0.8169 & 0.8549\\
\textbf{\texttt{expPROJ}} & 0.2176 & 0.6521 & \textbf{0.7224} & \textbf{0.7781} & \textbf{0.8225} & \textbf{0.8566}\\
\bottomrule
\end{tabular}
\caption{7B OlympiadBench. Both constraints exceed GRPO at large $k$ (\texttt{expPROJ} leads at every
$k\ge32$; pass@256 $0.8566$ vs $0.8497$), gap widening with $k$, at a small pass@1 cost.
Paired-bootstrap: pass@1 significant (\texttt{expPROJ} $-0.0178\,[-0.0209,-0.0148]$); pass@256 gain
\emph{within noise} ($+0.0070\,[-0.0122,+0.0262]$).}
\label{tab:oly7b}
\end{table}

\begin{table}[h]\centering\small
\begin{tabular}{lrrrrr}
\toprule
Arm (7B, OlympiadBench \emph{fragile band}, $n{=}329$, $1024$ samples) & pass@1 & pass@128 & pass@256 & pass@512 & pass@1024\\
\midrule
GRPO (control)      & \textbf{0.1895} & \textbf{0.9634} & 0.9833 & 0.9935 & 0.9970\\
\texttt{expSR}      & 0.1765 & 0.9625 & \textbf{0.9852} & \textbf{0.9964} & \textbf{1.0000}\\
\texttt{expPROJ}    & 0.1739 & 0.9614 & 0.9840 & 0.9961 & \textbf{1.0000}\\
\bottomrule
\end{tabular}
\caption{Tightened significance run at $k{=}1024$ on the fragile band. Both constraints reach
\emph{100\%} coverage at $k{=}1024$ (GRPO $0.997$) and lead at every $k\ge256$. Direction is stable,
but the effect is small ($\sim{+}0.003$) and the CI edge touches $0$
(\texttt{expPROJ}@1024 $+0.0030\,[0.0000,+0.0091]$): the fragile band \emph{saturates} at high $k$,
leaving little headroom---not a noise limitation but a difficulty-regime one.}
\label{tab:frag1024}
\end{table}

\begin{table}[h]\centering\small
\begin{tabular}{lrrrrr}
\toprule
Arm (7B, AIME 24/25/26, $n{=}90$) & pass@1 & pass@16 & pass@64 & pass@128 & pass@256\\
\midrule
base                & 0.0560 & 0.2956 & 0.4409 & 0.5041 & 0.5556\\
GRPO                & \textbf{0.0643} & \textbf{0.2997} & 0.4345 & 0.5093 & \textbf{0.6000}\\
\texttt{expSR}      & 0.0550 & 0.2817 & 0.4174 & 0.4937 & 0.5667\\
\texttt{expPROJ}    & 0.0553 & 0.2818 & 0.4176 & 0.4948 & 0.5889\\
\bottomrule
\end{tabular}
\caption{7B AIME (boundary condition). \textbf{No crossover here}: GRPO $\ge$ base at every $k$, so
there is no lost coverage to recover, and the constraint is $\approx$no-op-to-slightly-negative---as
predicted where the fragile-band premise fails. $n{=}90$ also gives CI half-width ${\pm}0.10$, so this
slice is underpowered at large $k$ regardless. \citet{yue2025} report the AIME crossover is largest at
\emph{14B/32B} and requires $k$ up to $1024$; our $7$B/$k{=}256$ run is below that regime.}
\label{tab:aime7b}
\end{table}

\paragraph{1.5B (in-domain + OOD).} On MATH-500 both constraints reach pass@256 $0.950$ vs GRPO
$0.946$ (recovering ${\sim}50\%$ of the base$\to$GRPO deficit), transferring to OlympiadBench OOD
(\texttt{expPROJ} $+0.014$ at pass@256). Paired-bootstrap CIs place all large-$k$ gains within noise
at feasible $n$; the pass@1 cost is small and significant ($\le0.011$).

\paragraph{Mechanism check.} \texttt{expSR} held $\approx100\%$ of banked base-correct modes above
their floor throughout training; \texttt{expPROJ} oscillated $\approx0.88$--$1.0$; plain GRPO has no
floor and drives fragile modes below the $1/K$ sampling threshold---the causal link between
\S\ref{sec:mech} and the coverage gains.

\paragraph{(Recipe) Why our own models under-show the crossover.}
\label{sec:recipe}
LoRA ($r{=}32$ and $r{=}256$) plus light training ($500$ steps) \emph{compress} the crossover: GRPO
loses only $3$--$14$ problems, versus Oat-Zero's $27$ (Table~\ref{tab:oatcross}) and the paper's
larger figures. LoRA structurally caps how far the policy can drift, hence how much coverage it can
destroy. The phenomenon---and thus the method's target---requires full-parameter, heavily-trained RL.
This is a methods lesson, not a negative result, and it dictates the recipe for \S\ref{sec:round2}.

%====================================================================
\section{Consequence: coverage as the substrate for continued RL}
\label{sec:round2}
The stakes of coverage are made concrete by continued (multi-round) RL. \textbf{Hypothesis:} round-1
RL prunes modes that round-2 RL needs, so preserving them raises the continued-RL \emph{accuracy}
ceiling. \textbf{Design (fork at round 1):} $\pi_0\!\to\!\{\pi_1^{\text{ctrl}},\pi_1^{\text{ours}}\}$
where $\pi_1^{\text{ours}}$ applies the floor \emph{during} round 1 (a ratchet, while modes are
alive---not restoration afterward, which degenerates to NLL replay); then run \emph{identical}
unconstrained round-2 RL on both. A third arm replays $\pi_1$'s lost set at round 2 to separate
\emph{preservation} from \emph{restoration}. The pass@1 cost of the floor becomes an investment
(``spend $1.8$ pts of round-1 pass@1, recover more at round-2''). Result R3 (exploitability gate,
$66\%$) confirms the design is viable before committing full-FT compute.

\paragraph{(R4) The experiment lands: coverage preservation raises the large-$k$ continued-RL ceiling.}
We ran the full-parameter path. Two round-1 forks of Qwen2.5-Math-7B---$\pi_1^{\text{ctrl}}$ (plain
GRPO) and $\pi_1^{\text{ours}}$ (GRPO $+$ support floor)---were each carried to $400$ steps, then given
\emph{identical} unconstrained full-FT GRPO round-2 RL to $100$ steps, and evaluated on the
$329$-problem Olympiad fragile band at $1024$ samples (Table~\ref{tab:round2}). The result is a
\textbf{second crossover, now between the two continued-RL models}: continued RL from the plain fork
(R2$\leftarrow$ctrl) wins small $k$ (pass@1 $0.378$ vs $0.299$---sharpening), but continued RL from the
coverage-preserved fork (R2$\leftarrow$ours) \textbf{overtakes at $k\ge32$ and the gap widens
monotonically} ($+0.0037$ at $k{=}64 \to +0.0068$ at $k{=}256$, reaching $0.990$). Both dominate base
everywhere. So the round-1 floor does not merely preserve modes for their own sake: the preserved
coverage is \emph{convertible by a second RL round into a higher large-$k$ ceiling} than the plain
fork can reach---the round-1 pass@1 tax is recovered as round-2 coverage, exactly the investment the
design predicted.

\begin{table}[h]\centering\small
\begin{tabular}{lrrrrrr}
\toprule
Arm (7B, Olympiad fragile band, $n{=}329$, $1024$ samples) & pass@1 & pass@16 & pass@32 & pass@64 & pass@128 & pass@256\\
\midrule
base                              & 0.1728 & 0.7519 & 0.8524 & 0.9175 & 0.9546 & 0.9763\\
R2$\leftarrow$ctrl (plain GRPO fork)      & \textbf{0.3780} & \textbf{0.8456} & 0.9067 & 0.9456 & 0.9696 & 0.9832\\
\textbf{R2$\leftarrow$ours (floor fork)}  & 0.2986 & 0.8381 & \textbf{0.9069} & \textbf{0.9494} & \textbf{0.9754} & \textbf{0.9900}\\
\bottomrule
\end{tabular}
\caption{\textbf{Continued-RL ceiling.} After identical round-2 GRPO, the coverage-preserved fork leads
at every $k\ge32$ (crossover at $k\!\approx\!32$; gap $+0.0037\!\to\!+0.0068$ over $k{=}64\!\to\!256$),
while the plain fork wins small $k$. Directionally consistent with the round-1 base-vs-GRPO crossover,
now reproduced \emph{between two round-2 models}. Gaps are small ($\sim{+}0.005$) and paired-bootstrap
CIs / the round-2 replay (restoration) arm are pending; direction and monotone-in-$k$ widening are
stable.}
\label{tab:round2}
\end{table}

%====================================================================
\section{Mechanism at the reasoning-mode level}
\label{sec:modelevel}
The results above are answer/pass@k-level. We now measure the mechanism directly on \emph{reasoning
modes} --- strategy families, not exact traces --- via a held-out bank of base-correct solutions
(11{,}664 witnesses on the Olympiad fragile band, clustered by strategy; the floor model never
trained on these traces). For each witness we compute $\Delta = \log\pi_\theta(y\mid q) - \log\pi_0(y\mid q)$
(teacher-forced), the same quantity the certificate (\S\ref{sec:method}) constrains. Define a mode as
\emph{collapsed} if $\Delta < -10$ nats (driven far below the base's probability, toward the sampling
floor).

\paragraph{(R5) Recoverability and extinction.} Recoverability $R_K(M)=1-(1-p_\theta(M))^K$. On the
329-problem fragile band, identical continued RL from the plain fork drove \textbf{2} base-recoverable
problems to \emph{zero} recoverable mass (unrecoverable at any $K$); from the coverage-preserved fork,
\textbf{0}. The extinction is invisible in pass@256 (0.983 vs 0.990) but categorical in recoverable-mode count.

\paragraph{(R6) Mode-mass certificate --- floor preserves what plain RL collapses (held out).}
Table~\ref{tab:certificate}: on traces the floor never trained on, plain GRPO collapses \textbf{39\%}
of base reasoning modes ($\Delta<-10$) while the floor collapses \textbf{2.3\%} ($\approx$17$\times$);
paired, the floor preserves more mass on 77\% of witnesses. Grpo's collapse rate is
bank-independent ($\approx$39--40\% on both its own and held-out banks) --- plain RL destroys base
mode-mass regardless of which traces we measure.

\begin{table}[h]\centering\small
\begin{tabular}{lrrr}
\toprule
Arm ($\Delta=\log\pi_\theta-\log\pi_0$ on base mode-witnesses) & mean $\Delta$ & \% $\ge\log\alpha$ & \% collapsed ($\Delta<-10$) \\
\midrule
\multicolumn{4}{l}{\emph{Held-out bank (11{,}664 witnesses; floor never trained on these):}}\\
plain GRPO (continued, 100 steps)     & $-10.11$ & 22.6\% & \textbf{39.0\%}\\
\textbf{GRPO + floor (continued, 100 steps)} & $\mathbf{-1.68}$ & \textbf{46.2\%} & \textbf{2.3\%}\\
\midrule
\multicolumn{4}{l}{\emph{Interventional (from base, identical 150-step GRPO; floor = sole difference):}}\\
plain GRPO                             & $-1.96$ & 46.2\% & 4.1\%\\
\textbf{GRPO + floor}                  & $\mathbf{+1.38}$ & \textbf{85.0\%} & \textbf{0.0\%}\\
\bottomrule
\end{tabular}
\caption{\textbf{The floor causally preserves reasoning-mode mass that plain RL erodes.} Top:
held-out mode-mass certificate (\S\ref{sec:modelevel}), removing any bank-provenance confound. Bottom:
interventional --- two arms from base with \emph{identical} data/steps, the off-policy floor the only
difference; plain RL erodes mode mass (mean $-1.96$, 4.1\% collapse), the floor raises it (mean
$+1.38$, \textbf{0\%} collapse), a $+3.34$-nat paired gap (floor $>$ plain on 69\% of witnesses). This
is the causal core: identical RL, floor prevents the collapse.}
\label{tab:certificate}
\end{table}

\paragraph{(R7) The $K\!\cdot\!p\approx1$ boundary predicts extinction.} The mechanism
(\S\ref{sec:mech}) says a mode is support-blind once $K\!\cdot\!p_\theta(M)<1$. At deployment $K{=}1024$,
plain GRPO leaves exactly the \emph{same} 2 problems below $K\!\cdot\!p{=}1$ that R5 found extinct
(\#168, \#237; the floor leaves 0) --- one of them a mode the base solved $33/1024$ times that plain RL
drove to $0$ and the floor preserved. The support-blind boundary and the observed extinctions coincide.

\emph{Honesty on the interventional run:} the $\Delta$-vs-$K\!\cdot\!p$ gradient is flat on this
concentrated-$p$ fragile band (the floor helps uniformly, $\approx+3.3$ nats across all bins), and a
25-step run showed no effect at all (too short to cross the boundary). The sharp $K\!\cdot\!p{=}1$
transition is carried by the observational analysis (R7); the interventional run establishes the
causal \emph{preservation}.

\paragraph{(R8) Baselines --- the floor beats global KL, not just vanilla GRPO.}
Table~\ref{tab:baselines}: a matched 150-step GRPO$+$global-KL arm ($\beta{=}0.04$, the canonical
coverage baseline) under the identical setup. Global KL genuinely helps (collapse $4.1\%\!\to\!0.5\%$,
mean $\Delta$ $-1.96\!\to\!-0.73$) --- it is a real baseline, not a strawman --- but the floor
dominates it: only the floor has \emph{positive} mean $\Delta$ (it \emph{raises} base mode mass),
$85\%$ vs $57\%$ of modes above the $\alpha{=}0.5$ floor, and $0\%$ vs $0.5\%$ collapse; paired,
floor $-$ global-KL $= +2.11$ nats (floor $>$ KL on $60\%$ of witnesses). Global KL applies symmetric,
uniform pressure toward base; the floor is one-sided and off-policy on the exact fragile modes, so its
signal stays strong precisely where those modes live. A symmetric base-anchoring arm (PBA's stronger
replay variant / DPH-RL family) is the strongest baseline yet ($66.8\%$ of modes above the floor) but
remains net-negative and still collapses $1.2\%$ of modes; the one-sided floor still leads it by
$+1.66$ nats (floor $>$ it on $62\%$). UCPO remains to be added.

\begin{table}[h]\centering\small
\begin{tabular}{lrrr}
\toprule
Arm (150 steps from base; $\Delta$ on base mode-witnesses) & mean $\Delta$ & \% $\ge\log\alpha$ & \% collapsed \\
\midrule
plain GRPO                    & $-1.96$ & 46.2\% & 4.1\%\\
GRPO $+$ global KL ($\beta{=}0.04$) & $-0.73$ & 57.1\% & 0.5\%\\
GRPO $+$ base-anchor (PBA/DPH-RL-style, symmetric) & $-0.28$ & 66.8\% & 1.2\%\\
\textbf{GRPO $+$ floor (ours, one-sided)} & $\mathbf{+1.38}$ & \textbf{85.0\%} & \textbf{0.0\%}\\
\bottomrule
\end{tabular}
\caption{\textbf{Baselines (mode-mass, 1055 witnesses; identical 150-step setup).} Ordering
plain $<$ global-KL $<$ base-anchor $<$ \textbf{floor}. The symmetric base-anchor (PBA's stronger
replay variant / DPH-RL family) is the strongest baseline ($66.8\%$ above the floor) but stays
net-negative and still collapses $1.2\%$ of modes; only the \emph{one-sided} floor achieves
\emph{positive} mean $\Delta$ (it raises mass), $85\%$ above the floor, and $0\%$ collapse
(floor $-$ base-anchor $=+1.66$ nats, floor $>$ it on $62\%$; floor $-$ global-KL $=+2.11$, floor $-$
plain $=+3.34$). The distinction is causal: symmetric anchoring taxes and still leaks; the one-sided
floor preserves without the tax.}
\label{tab:baselines}
\end{table}

%====================================================================
\section{What is and is not established}
\textbf{Established.} (1) A mechanism result (Prop.~\ref{prop:prune}--\ref{prop:rescue}) bounding an
entire family of on-policy methods. (2) A qualitatively different object---coverage as an off-policy
feasibility constraint---with two working implementations and a direct diagnostic. (3) A
\emph{directionally consistent} coverage gain over GRPO across scale (1.5B, 7B) and datasets
(MATH-500, OlympiadBench), at small bounded pass@1 cost.

\noindent\textbf{Not yet established (honest).} A \emph{statistically significant} large-$k$ coverage
win. On MATH-500/Olympiad at $k\le256$ the benchmarks are near-saturated (little headroom); on AIME at
$7$B there is no crossover to recover; the $k{=}1024$ fragile-band gain is real in sign but tiny at
the saturation ceiling. The theory is stated as propositions with proof sketches, not full theorems.

\paragraph{Predicted resolving regime (next experiments).} Prop.~\ref{prop:rescue} predicts the
effect is largest where a wide, \emph{un-saturated} fragile band exists. Per \citet{yue2025}, that is
harder benchmarks (Omni-MATH-Rule, AIME) on \emph{larger} models (14B/32B) evaluated to $k{=}1024$.
We therefore run the grid $\{$7B, 14B, 32B$\}\times\{$Olympiad, Omni-MATH, AIME$\}$ at $k{=}1024$ and
report the \texttt{expPROJ}$-$GRPO coverage gain as a function of scale and benchmark hardness---a
scaling-law test of the mechanism.

\begin{thebibliography}{9}
\bibitem{yue2025} Y.~Yue et al. \emph{Does Reinforcement Learning Really Incentivize Reasoning
Capacity in LLMs Beyond the Base Model?} arXiv:2504.13837, 2025.
\end{thebibliography}

\end{document}
